\documentclass[12pt,a4paper]{article}

\usepackage[utf8]{inputenc}
\usepackage[T1]{fontenc}
\usepackage[brazil]{babel}
\usepackage{amsmath,amssymb,mathtools}
\usepackage{graphicx}
\usepackage{xcolor}
\usepackage{float}
\usepackage{booktabs}
\usepackage{geometry}
\usepackage{caption}
\usepackage{placeins}
\usepackage{array}
\usepackage{listings}
\usepackage{hyperref}
\usepackage{fancyhdr}
\geometry{left=3cm,right=2cm,top=3cm,bottom=2cm}
\setlength{\headheight}{15pt}
\hypersetup{colorlinks=true, linkcolor=black, urlcolor=blue, citecolor=black}

\pagestyle{fancy}
\fancyhf{}
\rhead{Métodos Numéricos}
\lhead{Primeira Avaliação}
\cfoot{\thepage}

\lstdefinestyle{pythonstyle}{
    language=Python,
    basicstyle=\ttfamily\small,
    keywordstyle=\color{blue!70!black}\bfseries,
    commentstyle=\color{green!40!black},
    stringstyle=\color{red!50!black},
    showstringspaces=false,
    breaklines=true,
    frame=single,
    columns=fullflexible,
    keepspaces=true
}

\newcommand{\arquivofig}[1]{%
    \IfFileExists{figuras/#1}{%
        \includegraphics[width=0.85\textwidth]{figuras/#1}%
    }{%
        \fbox{\parbox[c][6cm][c]{0.80\textwidth}{\centering Arquivo de figura não encontrado:\\ \texttt{\detokenize{#1}}}}%
    }%
}

\title{Resolução da Primeira Avaliação de Métodos Numéricos}
\author{Gustavo Bittar Gonçalves}
\date{}

\begin{document}

\maketitle
\thispagestyle{fancy}

\section{Questão 1 -- Equação governante com condutividade térmica constante}
Para uma chapa bidimensional, transiente, sem geração interna de calor e com propriedades constantes, a equação governante da condução térmica é
\begin{equation}
\rho c_p \frac{\partial T}{\partial t} = k\left(\frac{\partial^2 T}{\partial x^2} + \frac{\partial^2 T}{\partial y^2}\right).
\end{equation}
Dividindo a equação por $\rho c_p$ e definindo a difusividade térmica como
\begin{equation}
\alpha = \frac{k}{\rho c_p},
\end{equation}
obtém-se
\begin{equation}
\frac{\partial T}{\partial t} = \alpha\left(\frac{\partial^2 T}{\partial x^2} + \frac{\partial^2 T}{\partial y^2}\right).
\label{eq:governante_constante}
\end{equation}

As condições de contorno informadas no enunciado são:
\begin{align}
T(0,y,t) &= T_{\text{imp}}, && 0<y<L_y, \\
\left.\frac{\partial T}{\partial x}\right|_{x=L_x} &= 0, && 0<y<L_y, \\
k\left.\frac{\partial T}{\partial y}\right|_{y=0} &= h_b\bigl(T(x,0,t)-T_{\infty,b}\bigr), && 0<x<L_x,
\end{align}
\begin{equation}
-k\left.\frac{\partial T}{\partial y}\right|_{y=L_y} = h_t\bigl(T(x,L_y,t)-T_{\infty,t}\bigr) + \varepsilon \sigma \left(T(x,L_y,t)^4 - T_{\text{viz}}^4\right), \quad 0<x<L_x,
\end{equation}
em que $h_b$ é o coeficiente convectivo da borda inferior, $h_t$ o coeficiente convectivo da borda superior, $\varepsilon$ a emissividade e $\sigma$ a constante de Stefan-Boltzmann.

A condição inicial é
\begin{equation}
T(x,y,0)=300\,\text{K}, \qquad 0<x<L_x,\;0<y<L_y.
\end{equation}

\section{Questão 2 -- Equação governante com condutividade variável}
Quando a condutividade varia linearmente com a temperatura,
\begin{equation}
k(T)=aT+b,
\end{equation}
a forma geral da equação de energia é
\begin{equation}
\rho c_p \frac{\partial T}{\partial t}=\nabla\cdot\bigl(k(T)\nabla T\bigr).
\end{equation}
Em coordenadas cartesianas bidimensionais,
\begin{equation}
\rho c_p \frac{\partial T}{\partial t}=
\frac{\partial}{\partial x}\left[k(T)\frac{\partial T}{\partial x}\right]+
\frac{\partial}{\partial y}\left[k(T)\frac{\partial T}{\partial y}\right].
\label{eq:governante_variavel}
\end{equation}
Substituindo $k(T)=aT+b$,
\begin{equation}
\rho c_p \frac{\partial T}{\partial t}=
\frac{\partial}{\partial x}\left[(aT+b)\frac{\partial T}{\partial x}\right]+
\frac{\partial}{\partial y}\left[(aT+b)\frac{\partial T}{\partial y}\right].
\end{equation}

As condições de contorno mantêm a mesma estrutura da Questão 1, apenas substituindo $k$ por $k(T)$ nas fronteiras convectivas/radiativas:
\begin{align}
T(0,y,t) &= T_{\text{imp}}, \\
\left.\frac{\partial T}{\partial x}\right|_{x=L_x} &= 0, \\
k(T)\left.\frac{\partial T}{\partial y}\right|_{y=0} &= h_b\bigl(T(x,0,t)-T_{\infty,b}\bigr), \\
-k(T)\left.\frac{\partial T}{\partial y}\right|_{y=L_y} &= h_t\bigl(T(x,L_y,t)-T_{\infty,t}\bigr) + \varepsilon \sigma \left(T(x,L_y,t)^4 - T_{\text{viz}}^4\right),
\end{align}
com condição inicial
\begin{equation}
T(x,y,0)=300\,\text{K}.
\end{equation}

\section{Questão 3 -- Discretização por diferenças finitas e Euler explícito}
Considere a malha uniforme $(x_i,y_j)$, com
\begin{equation}
\Delta x = \frac{L_x}{N_x-1},
\qquad
\Delta y = \frac{L_y}{N_y-1},
\qquad
\Delta t = t^{n+1}-t^n.
\end{equation}
Seja $T_{i,j}^n$ a temperatura no nó $(i,j)$ no instante $t^n$.

Para os pontos internos, aplicando diferenças finitas centrais de segunda ordem nos termos espaciais e Euler explícito no tempo à equação \eqref{eq:governante_constante}, resulta
\begin{equation}
\frac{T_{i,j}^{n+1}-T_{i,j}^{n}}{\Delta t}=
\alpha\left[
\frac{T_{i+1,j}^{n}-2T_{i,j}^{n}+T_{i-1,j}^{n}}{\Delta x^2}
+
\frac{T_{i,j+1}^{n}-2T_{i,j}^{n}+T_{i,j-1}^{n}}{\Delta y^2}
\right].
\end{equation}
Logo,
\begin{equation}
T_{i,j}^{n+1}=T_{i,j}^{n}+
\alpha\Delta t\left[
\frac{T_{i+1,j}^{n}-2T_{i,j}^{n}+T_{i-1,j}^{n}}{\Delta x^2}
+
\frac{T_{i,j+1}^{n}-2T_{i,j}^{n}+T_{i,j-1}^{n}}{\Delta y^2}
\right].
\label{eq:interna}
\end{equation}

\subsection{Tratamento das condições de contorno}
\subsubsection*{Borda esquerda: temperatura imposta em $x=0$}
Nesta fronteira a condição é de Dirichlet:
\begin{equation}
T_{0,j}^{n+1}=T_{\text{imp}}.
\end{equation}

\subsubsection*{Borda direita: adiabática em $x=L_x$}
A condição
\begin{equation}
\left.\frac{\partial T}{\partial x}\right|_{x=L_x}=0
\end{equation}
é implementada com nó fantasma, resultando em
\begin{equation}
\left.\frac{\partial^2 T}{\partial x^2}\right|_{i=N_x-1,j}
\approx \frac{2\left(T_{N_x-2,j}^{n}-T_{N_x-1,j}^{n}\right)}{\Delta x^2}.
\end{equation}
Assim, na borda direita,
\begin{equation}
T_{N_x-1,j}^{n+1}=T_{N_x-1,j}^{n}+\alpha\Delta t\left[
\frac{2\left(T_{N_x-2,j}^{n}-T_{N_x-1,j}^{n}\right)}{\Delta x^2}
+
\frac{T_{N_x-1,j+1}^{n}-2T_{N_x-1,j}^{n}+T_{N_x-1,j-1}^{n}}{\Delta y^2}
\right].
\end{equation}

\subsubsection*{Borda inferior: convecção em $y=0$}
Da condição
\begin{equation}
k\left.\frac{\partial T}{\partial y}\right|_{y=0}=h_b\left(T-T_{\infty,b}\right),
\end{equation}
obtém-se, com eliminação do nó fantasma,
\begin{equation}
\left.\frac{\partial^2 T}{\partial y^2}\right|_{i,0}
\approx
\frac{2\left(T_{i,1}^{n}-T_{i,0}^{n}\right)}{\Delta y^2}
-
\frac{2h_b}{k\Delta y}\left(T_{i,0}^{n}-T_{\infty,b}\right).
\end{equation}
Portanto,
\begin{equation}
T_{i,0}^{n+1}=T_{i,0}^{n}+\alpha\Delta t\left[
\frac{T_{i+1,0}^{n}-2T_{i,0}^{n}+T_{i-1,0}^{n}}{\Delta x^2}
+
\frac{2\left(T_{i,1}^{n}-T_{i,0}^{n}\right)}{\Delta y^2}
-
\frac{2h_b}{k\Delta y}\left(T_{i,0}^{n}-T_{\infty,b}\right)
\right].
\end{equation}

\subsubsection*{Borda superior: convecção + radiação em $y=L_y$}
Na borda superior,
\begin{equation}
-k\left.\frac{\partial T}{\partial y}\right|_{y=L_y}=h_t\left(T-T_{\infty,t}\right)+\varepsilon\sigma\left(T^4-T_{\text{viz}}^4\right).
\end{equation}
Eliminando o nó fantasma, obtém-se
\begin{equation}
\left.\frac{\partial^2 T}{\partial y^2}\right|_{i,N_y-1}
\approx
\frac{2\left(T_{i,N_y-2}^{n}-T_{i,N_y-1}^{n}\right)}{\Delta y^2}
-
\frac{2}{k\Delta y}
\left[
 h_t\left(T_{i,N_y-1}^{n}-T_{\infty,t}\right)+\varepsilon\sigma\left(\left(T_{i,N_y-1}^{n}\right)^4-T_{\text{viz}}^4\right)
\right].
\end{equation}
Dessa forma,
\begin{equation}
\begin{split}
T_{i,N_y-1}^{n+1}=&T_{i,N_y-1}^{n}+\alpha\Delta t\left[
\frac{T_{i+1,N_y-1}^{n}-2T_{i,N_y-1}^{n}+T_{i-1,N_y-1}^{n}}{\Delta x^2}
+
\frac{2\left(T_{i,N_y-2}^{n}-T_{i,N_y-1}^{n}\right)}{\Delta y^2} -\\
&-
\frac{2}{k\Delta y}
\left(
 h_t\left(T_{i,N_y-1}^{n}-T_{\infty,t}\right)+\varepsilon\sigma\left(\left(T_{i,N_y-1}^{n}\right)^4-T_{\text{viz}}^4\right)
\right)
\right].
\end{split}
\end{equation}

\subsection{Critério de estabilidade do esquema explícito}
Para o método explícito de Euler em duas dimensões, um limite clássico de estabilidade é
\begin{equation}
\Delta t \leq \frac{1}{2\alpha\left(\frac{1}{\Delta x^2}+\frac{1}{\Delta y^2}\right)}.
\end{equation}
No programa, foi adotado um fator de segurança adicional para garantir estabilidade numérica.

\section{Questão 4 -- Algoritmo numérico}
O algoritmo para resolver o problema do item 3 pode ser descrito como segue:

\begin{enumerate}
    \item Definir as propriedades físicas do material e os parâmetros geométricos da chapa.
    \item Definir as condições de contorno e a condição inicial.
    \item Escolher a malha $(N_x,N_y)$ e calcular $\Delta x$, $\Delta y$ e $\Delta t$.
    \item Inicializar a matriz de temperatura com $T=300$ K em todo o domínio e impor $T=400$ K em $x=0$.
    \item Para cada passo temporal:
    \begin{enumerate}
        \item atualizar os nós internos pela Eq. \eqref{eq:interna};
        \item atualizar a borda direita com a condição adiabática;
        \item atualizar a borda inferior com convecção;
        \item atualizar a borda superior com convecção e radiação;
        \item atualizar os cantos combinando as condições adjacentes;
        \item reimpor a condição de Dirichlet em $x=0$.
    \end{enumerate}
    \item Armazenar os campos de temperatura nos tempos pedidos: $0$, $1800$, $3600$, $7200$, $10800$ e $14400$ s.
    \item Ao final de $4$ h, extrair os perfis pedidos e os históricos temporais nos pontos especificados.
    \item Repetir o procedimento para as diferentes malhas solicitadas.
\end{enumerate}

\section{Questão 5 -- Implementação computacional}
Foi utilizado um programa em Python para resolver o problema transiente. O código utiliza \texttt{NumPy} para a solução em CPU e tenta utilizar \texttt{CuPy/CUDA} quando disponível. O programa também salva automaticamente todas as figuras pedidas no enunciado.

Nesta implementação foram adotados os seguintes dados do problema:
\begin{center}
\begin{tabular}{>{\raggedright\arraybackslash}p{7.0cm}c}
\toprule
\textbf{Parâmetro} & \textbf{Valor} \\
\midrule
Comprimento da chapa em $x$ & $L_x=1{,}0$ m \\
Comprimento da chapa em $y$ & $L_y=1{,}0$ m \\
Espessura & $5{,}0\times 10^{-3}$ m \\
Condutividade térmica do alumínio & $k=237$ W/(m.K) \\
Massa específica & $\rho=2700$ kg/m$^3$ \\
Calor específico & $c_p=897$ J/(kg.K) \\
Temperatura imposta em $x=0$ & $400$ K \\
Temperatura inicial & $300$ K \\
Convecção em $y=0$ & $h=100$ W/(m$^2$.K), $T_\infty=350$ K \\
Convecção em $y=L_y$ & $h=200$ W/(m$^2$.K), $T_\infty=280$ K \\
Radiação em $y=L_y$ & $\varepsilon=0{,}8$, $T_{\text{viz}}=300$ K \\
Tempo final & $4$ h \\
\bottomrule
\end{tabular}
\end{center}

\section{Resultados solicitados no item 5.1}
As figuras abaixo foram organizadas conforme a nomenclatura mostrada nos arquivos gerados pelo programa. Caso as imagens estejam na pasta \texttt{figuras/}, o documento compilará exibindo os gráficos automaticamente.

\subsection{Campos de temperatura para a malha de referência}
\begin{figure}[H]
\centering
\arquivofig{campo_temperatura_t_00000.png}
\caption{Campo de temperatura em $t=0$ s para a malha $N_x=N_y=10$.}
\end{figure}

\begin{figure}[H]
\centering
\arquivofig{campo_temperatura_t_01800.png}
\caption{Campo de temperatura em $t=1800$ s para a malha $N_x=N_y=10$.}
\end{figure}

\begin{figure}[H]
\centering
\arquivofig{campo_temperatura_t_03600.png}
\caption{Campo de temperatura em $t=3600$ s para a malha $N_x=N_y=10$.}
\end{figure}

\begin{figure}[H]
\centering
\arquivofig{campo_temperatura_t_07200.png}
\caption{Campo de temperatura em $t=7200$ s para a malha $N_x=N_y=10$.}
\end{figure}

\begin{figure}[H]
\centering
\arquivofig{campo_temperatura_t_10800.png}
\caption{Campo de temperatura em $t=10800$ s para a malha $N_x=N_y=10$.}
\end{figure}

\begin{figure}[H]
\centering
\arquivofig{campo_temperatura_t_14400.png}
\caption{Campo de temperatura em $t=14400$ s para a malha $N_x=N_y=10$.}
\end{figure}


\subsection{Perfis espaciais em \texorpdfstring{$t=4$ h}{t=4 h}}
\begin{figure}[H]
\centering
\arquivofig{perfil_x_meio.png}
\caption{Perfis de temperatura em \texorpdfstring{$x=L_x/2$}{x=Lx/2} para diferentes malhas.}
\end{figure}

\begin{figure}[H]
\centering
\arquivofig{perfil_x_quarto.png}
\caption{Perfis de temperatura em \texorpdfstring{$x=L_x/4$}{x=Lx/4} para diferentes malhas.}
\end{figure}

\begin{figure}[H]
\centering
\arquivofig{perfil_y_meio.png}
\caption{Perfis de temperatura em \texorpdfstring{$y=L_y/2$}{y=Ly/2} para diferentes malhas.}
\end{figure}

\begin{figure}[H]
\centering
\arquivofig{perfil_y_tres_quartos.png}
\caption{Perfis de temperatura em \texorpdfstring{$y=3L_y/4$}{y=3Ly/4} para diferentes malhas.}
\end{figure}

\subsection{Históricos temporais}
\begin{figure}[H]
\centering
\arquivofig{historico_centro.png}
\caption{Evolução temporal da temperatura em \texorpdfstring{$x=L_x/2$, $y=L_y/2$}{x=Lx/2, y=Ly/2} para diferentes malhas.}
\end{figure}

\begin{figure}[H]
\centering
\arquivofig{historico_quarto.png}
\caption{Evolução temporal da temperatura em \texorpdfstring{$x=L_x/4$, $y=L_y/4$}{x=Lx/4, y=Ly/4} para diferentes malhas.}
\end{figure}

\section{Conclusão}
A formulação apresentada permite resolver a condução de calor bidimensional transiente com condições de contorno mistas, incluindo adiabática, convecção e convecção associada à radiação. A discretização por diferenças finitas centrais de segunda ordem no espaço e Euler explícito no tempo conduz a um algoritmo simples e direto, adequado para implementação computacional. O programa em Python automatiza a solução para várias malhas e gera todas as figuras exigidas pela avaliação, facilitando a análise de independência de malha e da evolução temporal do campo térmico.

\appendix
\appendix
\section{Código-fonte em Python}

\lstinputlisting[style=pythonstyle]{main.py}
\end{document}
